% Source: https://artofproblemsolving.com/community/c3h1977774_fractionsdecimals_into_other_bases
% Stable identity: aops-topic-1977774-q01
(a) To convert a fraction, repeatedly multiply by the new base and read off the integer parts:
\[
\dfrac{5}{8}\cdot2=1+\dfrac{1}{4},\qquad\dfrac{1}{4}\cdot2=0+\dfrac{1}{2},\qquad\dfrac{1}{2}\cdot2=1+0 .
\]
The digits are $1,0,1$ and the process stops, so $\dfrac{5}{8}=\boxed{0.101_2}$. (Check: $\dfrac{1}{2}+\dfrac{1}{8}=\dfrac{5}{8}$.)

(b) Let $x=0.\overline{12}_3$. Shifting two places in base $3$ multiplies by $3^2=9$: $9x=12.\overline{12}_3=5+x$, since $12_3=1\cdot3+2=5$. So $8x=5$ and $x=\boxed{\dfrac{5}{8}}$.

(Both parts describe the same number: $\dfrac{5}{8}=0.101_2=0.\overline{12}_3$.)
